\begin{proposition}[Finite Borel gauges and an arithmetic escape]
\label{prop:gap-kernel-borel-gauges}
Retain $U_m(a)$ and $V_m(a)$ from
Proposition~\ref{prop:gap-kernel-green}.  For $0\leq r\leq3$, put
\begin{equation}
 W_m^{(r)}(a)
 =\frac{U_m(a)}{(m!)^{3-r}(a+1)_m^r}
 =\binom{a+m}{m}^{3-r}V_m(a),
 \label{eq:borel-gauges}
\end{equation}
and let
\[
 \mathcal W_r(s,z)=
 \sum_{a\geq1}\sum_{m\geq0}W_m^{(r)}(a)s^az^m.
\]
The following statements hold.

\begin{enumerate}
\item If
  \[
   \mathcal U(s,z)=\sum_{a\geq1}\sum_{m\geq0}U_m(a)s^az^m,
   \qquad \Theta=\theta_s+\theta_z,
  \]
  then
  \begin{equation}
   \boxed{
   \bigl(1-zP(\Theta)+z^2(\Theta+1)^6\bigr)\mathcal U(s,z)
   =\frac{s}{1-s}.}
   \label{eq:unnormalized-bivariate-PDE}
  \end{equation}
  Thus the integral array has a fixed differential equation, although it
  has factorial order three in the gap variable.

\item Successive gauges differ by the no-carry binomial coefficient:
  \begin{equation}
   W_m^{(r+1)}(a)=\binom{a+m}{m}^{-1}W_m^{(r)}(a).
   \label{eq:borel-gauge-transition}
  \end{equation}
  If $p\geq5$ and $a,m\geq0$ satisfy $a+m<p$, every multiplier in
  \eqref{eq:borel-gauges}--\eqref{eq:borel-gauge-transition} is a
  $p$-adic unit.  Hence all four gauges have exactly the same zero relation
  modulo~$p$ on the strict triangle.

\item None of the four series $\mathcal W_r$ is globally bounded.  More
  precisely, for every prime $p\geq7$,
  \begin{equation}
   v_p\bigl(W_p^{(r)}(1)\bigr)=-3
   \qquad(0\leq r\leq3).
   \label{eq:borel-gauge-denominator}
  \end{equation}
  Consequently no fixed rescaling of $s,z$ and of the whole series makes
  its coefficients integral.  In particular, none of the $\mathcal W_r$
  is algebraic over $\mathbb Q(s,z)$ or an ordinary diagonal or constant
  term of a fixed rational function regular at the expansion point.

\item Let
  \[
   \lambda_n=\operatorname {lcm}(1,\ldots,n),\qquad \lambda_0=1,
  \]
  and define the arithmetic gauge
  \begin{equation}
   \Xi_m(a)=6\lambda_{a+m}^3V_m(a).
   \label{eq:lcm-gap-gauge}
  \end{equation}
  Then $\Xi_m(a)\in\mathbb Z$ for all integers $a\geq1,m\geq0$, and
  $|\Xi_m(a)|\leq C^{a+m}$ for one absolute constant $C>1$.  Thus
  \[
   \mathcal X(s,z)=\sum_{a\geq1}\sum_{m\geq0}\Xi_m(a)s^az^m
  \]
  is a globally bounded bivariate series with positive radii of
  convergence.  If $p\geq7$ and $a+m<p$, then
  \begin{equation}
   U_m(a)=0\pmod p
   \quad\Longleftrightarrow\quad
   \Xi_m(a)=0\pmod p.
   \label{eq:lcm-gap-unit-equivalence}
  \end{equation}
  The factor $\lambda_n$ is not hypergeometric and cannot be generated by
  a finite product of fixed shifted factorials and a rational prefactor.
  Hence this construction lies outside the finite Borel/Gamma family; it
  shows that the denominator obstruction in part~(3) is not universal.

\item The arithmetic escape in part~(4) is not holonomic.  More precisely,
  \[
   \mathcal X(s,z)\quad\hbox{is not D-finite over }\mathbb Q(s,z).
  \]
  Consequently it is not an ordinary diagonal of a fixed rational
  function, nor a period germ of a fixed finite-rank algebraic
  Gauss--Manin system.

  This failure persists quantitatively after reduction.  Put
  \[
   \mathcal P(p)=
   \#\{q^e: q\text{ prime},\ e\geq1,\ 2\leq q^e\leq p-1\}.
  \]
  If \(p\geq11\) and \(Q_p\in\mathbb F_p[X]\) satisfies
  \[
   Q_p(n)=6\lambda_n^3\qquad(0\leq n<p),
  \]
  then
  \begin{equation}
   \deg Q_p\geq p-\mathcal P(p)
   =p-O\!\left(\frac p{\log p}\right).
   \label{eq:lcm-gauge-interpolation-degree}
  \end{equation}
  Thus even the zeroth \(z\)-slice has asymptotically full direct
  polynomial complexity modulo~\(p\).
\end{enumerate}
\end{proposition}

\begin{proof}
Multiplying the recurrence
\[
 U_{m+1}(a)=P(a+m)U_m(a)-(a+m)^6U_{m-1}(a)
\]
by $s^az^{m+1}$ and summing over $a\geq1,m\geq0$ gives
\eqref{eq:unnormalized-bivariate-PDE}; the missing $m=0$ boundary is
$\sum_{a\geq1}s^a=s/(1-s)$.  On the slice $a=1$, the Casoratian formula
gives
\begin{equation}
 U_m(1)=((m+1)!)^3c_{m+1}.
 \label{eq:unnormalized-boundary-growth}
\end{equation}
The recurrence bounds both independent scalar solutions by $C_0^n$, while
\eqref{eq:unnormalized-boundary-growth} contains the displayed factorial
cube.  More explicitly, $c_1>c_0\geq0$, and
\[
 P(n)-n^3>(n+1)^3\qquad(n\geq1),
\]
so the recurrence inductively gives $c_{n+1}>c_n>0$.  It also gives
$c_{n+1}<34c_n$.  Thus $1\leq c_n\leq34^{n-1}$ for $n\geq1$, proving the
asserted exact factorial order.  Formula
\eqref{eq:borel-gauges} follows immediately from
$(a+1)_m/m!=\binom{a+m}{m}$, and gives
\eqref{eq:borel-gauge-transition}.  On $a+m<p$, all factorials involved
belong to $\mathbb Z_p^\times$, proving part~(2).

We next propagate the denominator of the companion solution.  Proposition
\ref{prop:gap-kernel-denominator-obstruction} gives
\[
 p^3c_p\equiv1=b_0\pmod p.
\]
For $0\leq t\leq p-2$, multiply the recurrence at index $p+t$ by~$p^3$
and reduce modulo~$p$.  Induction gives
\begin{equation}
 p^3c_{p+t}\equiv b_t\pmod p
 \qquad(0\leq t\leq p-1).
 \label{eq:companion-denominator-propagation}
\end{equation}
In particular, $p^3c_{p+1}\equiv b_1=5\pmod p$ for $p\geq7$.  Since
\[
 W_m^{(r)}(1)=(m+1)^{3-r}c_{m+1},
\]
equation~\eqref{eq:borel-gauge-denominator} follows.  A fixed rescaling
multiplies the coefficient of $s z^p$ by a fixed scalar times a fixed
$p$-th power; for all sufficiently large primes these are units.  Thus the
new denominator primes cannot be removed.  Eisenstein's theorem for
algebraic power series gives the algebraic obstruction.  A rational
function regular at the origin has coefficient denominators supported on
one finite set of primes, and diagonal or constant-term extraction does not
enlarge that set; this proves the remaining assertions in part~(3).

For part~(4), Ap\'ery integrality says
\[
 6\lambda_n^3c_n\in\mathbb Z,
 \qquad b_n\in\mathbb Z.
\]
Writing $n=a+m$ and using~\eqref{eq:green-casoratian}, we obtain
\begin{align*}
 \Xi_m(a)=a^3\bigl(&b_{a-1}(6\lambda_n^3c_n)
             -(6\lambda_n^3c_{a-1})b_n\bigr).
 \label{eq:lcm-gap-integrality}
\end{align*}
The first parenthesized quantity is integral.  The second is integral
because $\lambda_{a-1}\mid\lambda_n$ and
\[
 6\lambda_n^3c_{a-1}
 =\left(\frac{\lambda_n}{\lambda_{a-1}}\right)^3
  (6\lambda_{a-1}^3c_{a-1}).
\]
This proves integrality.  The elementary estimate
$\lambda_n\leq4^n$, together with the recurrence bound
$|b_n|+|c_n|\leq C_0^n$, proves exponential growth in $a+m$.

On the strict triangle, both $6\lambda_{a+m}^3$ and $(a+1)_m^3$ are
$p$-adic units.  Comparing~\eqref{eq:lcm-gap-gauge} with
$U_m(a)=(a+1)_m^3V_m(a)$ proves
\eqref{eq:lcm-gap-unit-equivalence}.  Finally,
\[
 \frac{\lambda_{n+1}}{\lambda_n}
 =\begin{cases}
   q,&n+1\text{ is a positive power of the prime }q,\\
   1,&\text{otherwise}.
  \end{cases}
\]
If this ratio were a rational function of~$n$, its equality to~$1$ at
infinitely many integers would force it to be identically~$1$, contrary to
the infinitely many prime-power spikes.  A finite product of fixed shifted
Gamma factors has rational consecutive ratio.  This proves the final
scope statement in part~(4).

It remains to prove part~(5).  Put $A_n=6\lambda_n^3$.  We first show
that $(A_n)$ is not P-recursive.  Suppose to the contrary that
\begin{equation}
 \sum_{j=0}^r R_j(n)A_{n+j}=0
 \qquad(n\geq n_0)
 \label{eq:lcm-hypothetical-recurrence}
\end{equation}
for polynomials $R_0,\ldots,R_r\in\mathbb Q[X]$, not all zero.  The case
$r=0$ is immediate, so assume $r\geq1$.

For each $1\leq i\leq r$, choose two distinct primes
$\ell_i,\ell_i'$, with all $2r$ primes distinct, and impose
\[
 n\equiv-i\pmod{\ell_i\ell_i'}.
\]
The Chinese remainder theorem supplies infinitely many such positive
integers~$n$.  Every $n+i$ is divisible by two distinct primes and is
therefore not a prime power.  For all sufficiently large members of the
progression,
\[
 A_n=A_{n+1}=\cdots=A_{n+r}.
\]
Substitution in~\eqref{eq:lcm-hypothetical-recurrence} shows that
\[
 S(X):=\sum_{j=0}^rR_j(X)
\]
vanishes at infinitely many integers, and hence $S=0$.

Fix $1\leq k\leq r$.  For each nonzero integer
$d\in\{1-k,\ldots,r-k\}$, choose two new distinct primes
$\ell_d,\ell_d'>r+1$, all mutually distinct, and impose
\[
 q\equiv-d\pmod{\ell_d\ell_d'}.
\]
The resulting residue class is coprime to its modulus.  Dirichlet's theorem
therefore gives infinitely many primes~$q$ in that class.  Put $n=q-k$.
Every $q+d$ in the recurrence window, except $q$ itself, is divisible
by two distinct primes and is not a prime power.  Since
\[
 \frac{A_t}{A_{t-1}}=
 \begin{cases}
  q^3,&t=q^e\text{ for a prime }q,\\
  1,&t\text{ is not a prime power},
 \end{cases}
\]
there is a nonzero common value $B$ such that
\[
 A_{n+j}=B\quad(j<k),\qquad
 A_{n+j}=q^3B\quad(j\geq k).
\]
Using $S=0$, equation~\eqref{eq:lcm-hypothetical-recurrence} becomes
\[
 (q^3-1)\sum_{j=k}^rR_j(q-k)=0.
\]
It holds for infinitely many primes~$q$, so
$\sum_{j=k}^rR_j=0$ as a polynomial.  Descending from $k=r$ to~$1$
gives $R_r=\cdots=R_1=0$, and then $S=0$ gives $R_0=0$, a
contradiction.  Hence $(A_n)$ is not P-recursive.

The zeroth $z$-slice of the arithmetic gauge is
\[
 [z^0]\mathcal X(s,z)=6\sum_{a\geq1}\lambda_a^3s^a.
\]
Coefficient extraction from a multivariate D-finite series preserves
D-finiteness, whereas a univariate series is D-finite precisely when its
coefficient sequence is P-recursive.  Thus $\mathcal X$ is not D-finite.
Ordinary rational diagonals and finite-rank algebraic Gauss--Manin period
germs are D-finite, giving the two stated consequences.

Finally, let $D_p(X)=Q_p(X)-Q_p(X-1)$.  If
$1\leq n<p$ is not a prime power, then
$\lambda_n=\lambda_{n-1}$, so $D_p(n)=0$.  On the other hand,
\[
 D_p(2)=6(\lambda_2^3-\lambda_1^3)=42\ne0\pmod p
 \qquad(p\geq11).
\]
Thus $D_p$ is nonzero and has at least
$(p-1)-\mathcal P(p)$ distinct roots.  Since
$\deg D_p\leq\deg Q_p-1$, this proves the first inequality in
\eqref{eq:lcm-gauge-interpolation-degree}.  The number of prime powers of
exponent at least two below~$p$ is $O(\sqrt p\log p)$, while
$\pi(p)=O(p/\log p)$; the displayed asymptotic follows.
\end{proof}

\begin{remark}[What the lcm gauge does not yet provide]
\label{rem:lcm-gap-gauge-scope}
The series $\mathcal X$ is an integral positive-radius encoding of the exact
nonwrapping zero relation.  Proposition~\ref{prop:gap-kernel-borel-gauges}
shows that it is not D-finite and that direct polynomial models of its
reductions have degree $p-O(p/\log p)$.  It does not decide differential
algebraicity, and it does not exclude every nonlinear or
characteristic-dependent trace realization.  More importantly, on the
strict triangle the gauge only multiplies every gap equation by a unit, so
all far-bridge incidences and selected short-chain windows are exactly
unchanged.  The lcm gauge therefore identifies a genuine escape from the
denominator obstruction, but not an incidence estimate.
\end{remark}
